\scanpage{9}
\mspage{7}
\unsure{en} associant à \unsure{tt} chemin de $x$ dans $y$ le $\Sigma$-homomorphisme de $x$ \unsure{dans} $y$ obtenu en \uline{inversant} \deleted{les} (dans $X\Sigma^{-1}$) les \deleted{flèches} \inserted{tronçons} rétrogrades (ce qui est possible, \unsure{puisqu'ils|puisqu'elles} sont dans $\Sigma$), de sorte qu'on trouve dans $X\Sigma^{-1}$ un ``chemin direct''
\diagram{\begin{tikzcd}[ampersand replacement=\&,column sep=.65em,scale=.82,transform shape]
x=A_{-n}\arrow[r] \& \cdots\deleted{A}\arrow[r] \& A_{-1}\arrow[r] \& A_0\arrow[r] \& A_1\arrow[r,dash,"{\unsure{\text{---}|\to}}" description] \& A_2 \& \cdots \& A_n=A_{n+1}=\cdots=y\,,
\end{tikzcd}}
et on compose tous ces morphismes.

En fait, on trouve \unsure{aussitôt|aussi tôt} que \deleted{\unsure{pour}} l'on \deleted{\unsure{a}} \inserted{définit ainsi} un \uline{foncteur}
\[
\underline{\mathrm{Ch}}_\Sigma(X)\longrightarrow\underline{\mathrm{Fl}}(X\Sigma^{-1}),
\]
du fait que si on a un diagramme
\diagram{\begin{tikzcd}[ampersand replacement=\&]
A_{i-1}\arrow[d,"f_{i-1}"] \& A_i\arrow[l,"u_i"']\arrow[d,"f_i"] \& (i\text{ pair})\\
A'_{i-1} \& A'_i\arrow[l,"u'_i"']
\end{tikzcd}}
avec $u_i,u'_i\deleted{\unsure{-}}$ dans $\Sigma$, ce diagramme, en tant que diagramme dans $X\Sigma^{-1}$, est commutatif avec flèches horizontales des \uline{iso}, et \unsure{reste} donc commutatif en inversant les dites flèches. D'où un foncteur induit
\[
\underline{\mathrm{Ch}}_\Sigma(X;\unsure{x},y)\longrightarrow\underline{\mathrm{Fl}}(\underset{\inserted{X'}}{\underbrace{X\Sigma^{-1}}};x,y).
\]
Mais \deleted{la catégorie $\underline{\mathrm{Fl}}$} pour tte catégorie $X'$, et des objets $x,y$ de $X'$, la catégorie $\underline{\mathrm{Fl}}(X';x,y)$ est

% archive sheet number: 8
